% Slug: visual-gc-transition
% Source: figures/gc-color-transition.tex:3 (tikzpicture env start line)
% Original label: \label{fig:visual-gc-transition}
% Original caption (LaTeX source):
%   Marking moves the scanned object to black and its newly discovered children to gray without discarding other pending roots.
% Original caption (plain text):
%   Marking moves the scanned object to black and its newly discovered children to gray without discarding other pending roots.
% Required TikZ setup: \usepackage{tikz} + \usetikzlibrary{arrows.meta, positioning, shapes.geometric, fit} (loaded by support/compiler-pseudocode.sty).
% Shared styles used in this picture: `basic block', `cfg edge'.
%   Defined in support/compiler-pseudocode.sty: `basic block', `cfg edge', `back edge', `block label'.
% Usage: % Slug: visual-gc-transition
% Source: figures/gc-color-transition.tex:3 (tikzpicture env start line)
% Original label: \label{fig:visual-gc-transition}
% Original caption (LaTeX source):
%   Marking moves the scanned object to black and its newly discovered children to gray without discarding other pending roots.
% Original caption (plain text):
%   Marking moves the scanned object to black and its newly discovered children to gray without discarding other pending roots.
% Required TikZ setup: \usepackage{tikz} + \usetikzlibrary{arrows.meta, positioning, shapes.geometric, fit} (loaded by support/compiler-pseudocode.sty).
% Shared styles used in this picture: `basic block', `cfg edge'.
%   Defined in support/compiler-pseudocode.sty: `basic block', `cfg edge', `back edge', `block label'.
% Usage: % Slug: visual-gc-transition
% Source: figures/gc-color-transition.tex:3 (tikzpicture env start line)
% Original label: \label{fig:visual-gc-transition}
% Original caption (LaTeX source):
%   Marking moves the scanned object to black and its newly discovered children to gray without discarding other pending roots.
% Original caption (plain text):
%   Marking moves the scanned object to black and its newly discovered children to gray without discarding other pending roots.
% Required TikZ setup: \usepackage{tikz} + \usetikzlibrary{arrows.meta, positioning, shapes.geometric, fit} (loaded by support/compiler-pseudocode.sty).
% Shared styles used in this picture: `basic block', `cfg edge'.
%   Defined in support/compiler-pseudocode.sty: `basic block', `cfg edge', `back edge', `block label'.
% Usage: % Slug: visual-gc-transition
% Source: figures/gc-color-transition.tex:3 (tikzpicture env start line)
% Original label: \label{fig:visual-gc-transition}
% Original caption (LaTeX source):
%   Marking moves the scanned object to black and its newly discovered children to gray without discarding other pending roots.
% Original caption (plain text):
%   Marking moves the scanned object to black and its newly discovered children to gray without discarding other pending roots.
% Required TikZ setup: \usepackage{tikz} + \usetikzlibrary{arrows.meta, positioning, shapes.geometric, fit} (loaded by support/compiler-pseudocode.sty).
% Shared styles used in this picture: `basic block', `cfg edge'.
%   Defined in support/compiler-pseudocode.sty: `basic block', `cfg edge', `back edge', `block label'.
% Usage: \input{graphs-code/visual-gc-transition.tikz.tex} inside a figure environment
%   (wrap with \begin{figure} ... \caption{...} \label{fig:visual-gc-transition} ... \end{figure}).
% Note: tikzpicture body below is verbatim from the source; do not reformat.
\begin{tikzpicture}[node distance=11mm and 15mm]
    \node[basic block,align=center] (before)
      {Before scanning A\\white: B, C, E, F\\gray: A, D\\black: empty};
    \node[basic block,align=center,right=25mm of before] (after)
      {After scanning A\\white: E, F\\gray: D, B, C\\black: A};
    \draw[cfg edge] (before) -- node[above,font=\small] {scan A} (after);
    \node[font=\sffamily\small,align=center,below=of after] (note)
      {D remains pending.\\Scanning A discovers B and C.};
    \draw[cfg edge,dashed] (note) -- (after);
  \end{tikzpicture}
 inside a figure environment
%   (wrap with \begin{figure} ... \caption{...} \label{fig:visual-gc-transition} ... \end{figure}).
% Note: tikzpicture body below is verbatim from the source; do not reformat.
\begin{tikzpicture}[node distance=11mm and 15mm]
    \node[basic block,align=center] (before)
      {Before scanning A\\white: B, C, E, F\\gray: A, D\\black: empty};
    \node[basic block,align=center,right=25mm of before] (after)
      {After scanning A\\white: E, F\\gray: D, B, C\\black: A};
    \draw[cfg edge] (before) -- node[above,font=\small] {scan A} (after);
    \node[font=\sffamily\small,align=center,below=of after] (note)
      {D remains pending.\\Scanning A discovers B and C.};
    \draw[cfg edge,dashed] (note) -- (after);
  \end{tikzpicture}
 inside a figure environment
%   (wrap with \begin{figure} ... \caption{...} \label{fig:visual-gc-transition} ... \end{figure}).
% Note: tikzpicture body below is verbatim from the source; do not reformat.
\begin{tikzpicture}[node distance=11mm and 15mm]
    \node[basic block,align=center] (before)
      {Before scanning A\\white: B, C, E, F\\gray: A, D\\black: empty};
    \node[basic block,align=center,right=25mm of before] (after)
      {After scanning A\\white: E, F\\gray: D, B, C\\black: A};
    \draw[cfg edge] (before) -- node[above,font=\small] {scan A} (after);
    \node[font=\sffamily\small,align=center,below=of after] (note)
      {D remains pending.\\Scanning A discovers B and C.};
    \draw[cfg edge,dashed] (note) -- (after);
  \end{tikzpicture}
 inside a figure environment
%   (wrap with \begin{figure} ... \caption{...} \label{fig:visual-gc-transition} ... \end{figure}).
% Note: tikzpicture body below is verbatim from the source; do not reformat.
\begin{tikzpicture}[node distance=11mm and 15mm]
    \node[basic block,align=center] (before)
      {Before scanning A\\white: B, C, E, F\\gray: A, D\\black: empty};
    \node[basic block,align=center,right=25mm of before] (after)
      {After scanning A\\white: E, F\\gray: D, B, C\\black: A};
    \draw[cfg edge] (before) -- node[above,font=\small] {scan A} (after);
    \node[font=\sffamily\small,align=center,below=of after] (note)
      {D remains pending.\\Scanning A discovers B and C.};
    \draw[cfg edge,dashed] (note) -- (after);
  \end{tikzpicture}
